\begin{question}{30}{
    Tijdens een NAVO-oefening onderzoekt de Luchtmacht of twee verschillende typen boordcomputers (type A en type B) leiden tot verschillen in de reactietijd van piloten tijdens een standaard noodprocedure.
    Van beide typen boordcomputers worden onafhankelijke steekproeven genomen, waarbij de reactietijd wordt gemeten in seconden.
    \begin{center}
        \begin{tabular}{c|ccccccccccc}
            \toprule
                \textbf{Reactietijd type $A$ (in s)} & $4.8$ & $5.1$ & $4.5$ & $5.4$ & $4.9$ & $5.0$ & $4.7$ & $5.3$ & $4.6$ & $5.2$ & \\
                \textbf{Reactietijd type $B$ (in s)} & $4.9$ & $5.1$ & $4.4$ & $5.0$ & $5.7$ & $5.3$ & $4.8$ & $5.2$ & $5.5$ & $5.9$ & $4.2$ \\
            \bottomrule
        \end{tabular}
    \end{center}
    De Luchtmacht wil graag weten of de gemiddelde reactietijd van haar piloten significant verschilt bij gebruik van de twee typen boordcomputers.
    Je mag aannemen dat de reactietijden in beide gevallen normaal verdeeld zijn met onbekende verwachtingswaarden en standaardafwijkingen.
}   
    
    \subquestion{8}{
        Bereken voor beide steekproeven het steekproefgemiddelde en de steekproefvariantie.
    }

    \solution{
        We berekenen het steekproefgemiddelde $\overline{x}$ van de reactietijden met de boordcomputer van type $A$ als volgt:
        \begin{align*}
            \overline{x} = \frac{x_1+x_2+\ldots+x_n}{n} = \frac{ 4.8 + 5.1 + \ldots + 5.2 }{ 10 } = 4.95.\rubric{2}
        \end{align*}

        Dit steekproefgemiddelde kunnen we gebruiken om de steekproefvariantie $s_A^2$ te berekenen:
        \begin{align*}
            s_A^2 &= \frac{ (x_1 - \overline{x})^2 + (x_2 - \overline{x})^2 + \ldots + (x_n - \overline{x})^2 }{ n - 1 } \\
                &= \frac{ (4.8 - 4.95)^2 + (5.1 - 4.95)^2 + \ldots + (5.2 - 4.95)^2 }{ 10 - 1 } \\
                &\approx 0.0917. \rubric{2}
        \end{align*}

        Op eenzelfde manier berekenen het steekproefgemiddelde $\overline{y}$ van de reactietijden met de boordcomputer van type $B$ als volgt
        \begin{align*}
            \overline{y} = \frac{y_1+y_2+\ldots+y_m}{m} = \frac{ 4.9 + 5.1 + \ldots + 4.2 }{ 11 } \approx 5.0909.\rubric{2}
        \end{align*}

        We berekenen de steekproefvariantie $s_B^2$ als volgt:
        \begin{align*}
            s_B^2 &= \frac{ (y_1 - \overline{y})^2 + (y_2 - \overline{y})^2 + \ldots + (y_n - \overline{y})^2 }{ m - 1 } \\
                &= \frac{ (4.9 - 5.0909)^2 + (5.1 - 5.0909)^2 + \ldots + (4.2 - 5.0909)^2 }{ 11 - 1 } \\
                &\approx 0.2649.\rubric{2}
        \end{align*}
    }

    \subquestion{10}{
        Voer een $F$-toets uit om te bepalen of de varianties in de reactietijden gelijk kunnen zijn voor beide soorten boordcomputers.
        Kies voor het significantieniveau $\alpha = 0.05$.
    }
    \solution{
        In de vraag staat dat we aan mogen nemen dat $X_A \sim N(\mu_A = \text{ ? }, \sigma_A = \text{ ? })$ en $X_B \sim N(\mu_B = \text{ ? }, \sigma_B = \text{ ? })$.

        Bij een $F$-toets we om de varianties $\sigma_A^2$ en $\sigma_B$ gelijk zouden kunnen zijn, oftewel
        \begin{align*}
            H_0: \quad \sigma_A^2 = \sigma_B^2 \\
            H_1: \quad \sigma_A^2 \neq \sigma_B^2 \rubric{2}
        \end{align*}

        Verder is gegeven dat we mogen werken met een significantieniveau $\alpha=0.05$, en steekproefdata is al verzameld voor beide populaties.
        De toetsingsgrootheid voor een $F$-toets is gelijk aan
        \[
            F = \frac{S_A^2}{S_B^2}, 
        \]
        en volgt onder $H_0$ een $F$-verdeling met $\dfn = n-1 = 9$ en $\dfd = m-1 = 10$ vrijheidsgraden.\rubric{2}

        De geobserveerde toetsingsgrootheid is gelijk aan
        \[
            f = \frac{s_A^2}{s_B^2} = \frac{0.0917}{0.2649} \approx 0.3460. \rubric{2}
        \]
        
        We voeren een tweezijdige toets uit (dit is altijd het geval bij een $F$-toets voor gelijke varianties), oftewel de $p$-waarde is gelijk aan het minimum van de linker- en de rechteroverschrijdingskans van $f$.
        In dit geval is de linkeroverschrijdingskans kleiner:
        \begin{align*}
            p = P(F \ge f) = \text{Fcdf}(\text{lower}=0, \text{upper}=0.3460, \text{df1}=9, \text{df2}=10) \approx 0.0628.\rubric{2}
        \end{align*}
        Aangezien de $p$-waarde groter is dan $\alpha/2$ (tweezijdige toets!), wordt de nulhypothese $H_0$ aangenomen.\rubric{1}
        Er is onvoldoende bewijs om op basis van deze steekproef de aanname van gelijke varianties te verwerpen. \rubric{1}
    }

    \subquestion{12}{
        Bepaal met behulp van een onafhankelijke $t$-toets of de gemiddelde reactietijd significant verschil bij gebruik van de twee verschillende typen boordcomputers.
        Kies opnieuw voor het significantieniveau $\alpha = 0.05$.
    }
    \solution{
        We toetsen of de gemiddelde reactietijden $\mu_A$ en $\mu_B$ voor beide typen boordcomputers significant van elkaar verschilt.
        De hypotheses kunnen we daarom als volgt defini\"eren:
        \begin{align*}
            H_0: &\quad \mu_A = \mu_B \text{ (geen significant verschil) } \\
            H_1: &\quad \mu_A \neq \mu_B \text{ (wel een significant verschil) } \rubric{2}
        \end{align*}

        Verder is gegeven dat we mogen werken met een significantieniveau $\alpha=0.05$, en steekproefdata is al verzameld voor beide populaties.

        Op basis van ons antwoord bij vraag (b) mogen we aannemen dat $\sigma = \sigma_A = \sigma_B$.
        Dat betekent dat we met de \enquote{pooled variance} mogen werken als schatting voor de gemeenschappelijke onbekende $\sigma$. \rubric{1}
        
        \begin{align*}
            s_P^2 = \frac{(n-1)\cdot s_A^2 + (m-1) \cdot s_B^2}{n-1+m-1} = \frac{9 \cdot 0.0917 + 10\cdot 0.2649}{19} \approx 0.1828. \rubric{3}
        \end{align*}
        
        De toetsingsgrootheid van de bijbehorende $t$-toets is gelijk aan
        \[
            t = \frac{(\overline{x}-\overline{y}) - (\mu_A - \mu_B)}{\sqrt{\frac{s_P^2}{n} + \frac{s_P^2}{m}}} \approx -0.7542, 
        \]
        en komt uit een $t$-verdeling met $\text{df}=19$ vrijheidsgraden. \rubric{2}
        Omdat we tweezijdig toetsen, is het kritieke gebied van de vorm $(-\infty; g_1]$ en $[g_2, \infty)$.
        Deze kritieke grenzen kunnen we met de $t$-verdeling bepalen, namelijk
        \begin{align*}
            g_1 &= \invt(\text{opp} = \alpha/2, \text{df}=n+m-2) = \invt(\text{opp}=0.025, \text{df}=19) \approx -2.0930 \\ 
            g_2 &= \invt(\text{opp} = 1-\alpha/2, \text{df}=n+m-2) = \invt(\text{opp}=0.975, \text{df}=19) \approx 2.0930 \rubric{2}
        \end{align*}
        
        De toetsingsgrootheid $t\approx -0.7542$ ligt dus niet in het kritieke gebied, dus $H_0$ wordt in dit geval aangenomen.\rubric{1}
        Er is op basis van deze steekproeven onvoldoende reden om aan te nemen dat de gemiddelde reactietijd significant verschilt bij gebruik van de twee typen boordcomputers.
        \rubric{1}
    }
\end{question}